\begin{theorem}[单参对数读出与对角多通道读出的分辨率分岔：常数精度的必要代价]\label{thm:spg-single-parameter-log-readout-vs-diagonal-multichannel-bifurcation}
取互异素数 \(p_1,\dots,p_k\) 并令 \(b\in\{0,1\}^k\) 表示 \(k\)-比特 Gödel 寄存器。
定义单参对数读出
\[
u(b):=\sum_{i=1}^k b_i\log p_i.
\]
则存在 \(b\neq b'\) 使得
\[
\abs{u(b)-u(b')}
\ \le\
\frac{1}{2^k-1}\sum_{i=1}^k \log p_i.
\]
因此，任何仅观测单标量 \(u(b)\) 且具有误差尺度不小于上述最小间距量级的机制，都不可能对全部 \(b\in\{0,1\}^k\) 保证唯一译码；若要求无碰撞译码，则精度预算必随 \(k\) 指数增长。
\par
相对地，在推论 \ref{cor:spg-godel-axis-boundary-invertibility} 与定理 \ref{thm:spg-godel-boundary-decoding-relative-error-threshold} 的设定下，若可获得对角多通道比值
\[
r_i:=\frac{dB_{ii}}{B_0}=p_i^{2b_i},
\]
并满足相对误差模型
\(
\widetilde B_0=B_0(1+\delta_0),\ 
\widetilde B_{ii}=B_{ii}(1+\delta_i)
\)
且 \(|\delta_0|,|\delta_i|\le\varepsilon\)，则当
\[
\varepsilon<\frac{p_{\min}-1}{p_{\min}+1},
\qquad
p_{\min}:=\min_i p_i,
\]
最近整数译码
\[
\widetilde b_i:=\left\lfloor \frac{1}{2\log p_i}\log\!\left(\frac{d\,\widetilde B_{ii}}{\widetilde B_0}\right)\right\rceil
\]
可保证对所有 \(i\) 精确恢复 \(b_i\)。
从而，在“单参坐标”与“多通道对角坐标”之间存在严格分岔：前者的无碰撞译码要求指数级精度，后者可在与 \(k\) 无关的常数相对误差阈值下实现逐坐标精确译码。
\leanverified{paper\_spg\_single\_parameter\_log\_readout\_vs\_diagonal\_multichannel\_bifurcation}
\end{theorem}
\begin{proof}
把集合 \(\{u(b):b\in\{0,1\}^k\}\) 视为实轴上的 \(2^k\) 个点，且
\(
0\le u(b)\le \sum_{i=1}^k\log p_i
\)。
由鸽巢原理，区间 \([0,\sum_i\log p_i]\) 被这 \(2^k\) 个点划分出的相邻间距最小值不超过区间长度除以 \(2^k-1\)，从而存在 \(b\neq b'\) 满足所示不等式。
对单参读出的不可唯一性结论由“误差大于最小间距则发生碰撞”直接推出。
\par
多通道情形由定理 \ref{thm:spg-godel-boundary-decoding-relative-error-threshold} 直接给出。证毕。
\leanpartial{paper\_spg\_single\_parameter\_log\_readout\_pigeonhole\_k1}{仅 $k=1$ 实例 + Gödel 推论；通用 $k$ 鸽巢与论文 (B) 多通道译码待后轮}
\leanverified{weightedSum}
\leanverified{pigeonhole\_k\_one}
\leanverified{godel\_log\_readout\_pigeonhole\_k1}
\end{proof}

\par
上述单参对数读出的鸽巢论证可抽象为：在单标量投影下，离散态集的最小间距由其取值区间长度与点数之比控制。
对一般加权能量壳上的回路格点，同一机制给出由回路圆维驱动的幂律碰撞阈值，并可与附录的椭球格点计数定理严格对齐。

\begin{theorem}[Stokes--Gödel 椭球投影的最小间距律：单椭圆门的必然碰撞阈值]\label{thm:spg-stokes-godel-ellipsoid-min-spacing-single-gate}
令 \(G=(V,E)\) 为有限连通无自环多重图，并记回路圆维
\[
r:=\operatorname{rank}_{\ZZ}H_1(G;\ZZ)=|E|-|V|+1.
\]
给定正权 \(w:E\to(0,\infty)\)，令能量二次型为
\(
\mathcal E_w(\varphi):=\sum_{e\in E}w_e\varphi_e^2
\)。
取回路格 \(\Lambda(G)=\ker(\partial)\subset \ZZ^{E}\) 的任一 \(\ZZ\)-基坐标 \(k\in\ZZ^r\)，并令 \(A\in\mathrm{Sym}_r^{+}(\RR)\) 为该能量在回路坐标上的 Gram 矩阵（例如由定理 \ref{thm:graph-energy-shell-lattice-counting} 的构造给出）。
对能量预算 \(E>0\) 定义椭球格点集
\[
\mathcal M_A(E):=\left\{k\in\ZZ^r:\ k^{\mathsf T}Ak\le E\right\}.
\]
取任意非零线性观测 \(\ell(k):=\langle a,k\rangle\)（\(a\in\RR^r\setminus\{0\}\)）。
在 Gödel--Stokes 字典下，\(\mathcal M_A(E)\) 可被视为能量预算内的回路可行域，而 \(\ell\) 可取为 Gödel 码长的对数线性泛函（例如取 \(a=\log p\)）。
则存在常数 \(C(a,A)>0\) 使得当 \(E\to\infty\) 时，必存在不同 \(k\neq k'\in\mathcal M_A(E)\) 满足
\[
\abs{\ell(k)-\ell(k')}
\le
C(a,A)\,E^{\frac{1-r}{2}}+O\!\left(E^{-\frac r2}\right),
\]
其中可取显式
\[
C(a,A)=\frac{2\sqrt{\det A}}{\operatorname{Vol}(B_r)}\,
\sqrt{a^{\mathsf T}A^{-1}a}.
\]
特别地，当 \(r\ge 2\) 时最小间距呈 \(\asymp E^{(1-r)/2}\) 的衰减；因此任何仅通过单标量 \(\ell\)（等价地：单个对数尺度参数）进行的读出，一旦精度不优于该量级，则在能量预算 \(E\) 内不可避免地产生碰撞。
\leanverified{paper\_spg\_stokes\_godel\_ellipsoid\_min\_spacing\_single\_gate}
\end{theorem}
\begin{proof}
由定理 \ref{thm:graph-energy-shell-lattice-counting}（取 \(\delta=0\)）的体积主项，有
\[
|\mathcal M_A(E)|
=
\frac{\operatorname{Vol}(B_r)}{\sqrt{\det A}}\,E^{r/2}
+O\!\left(E^{(r-1)/2}\right).
\]
另一方面，由命题 \ref{prop:graph-energy-ellipsoid-godel-codelength-dual} 的对偶极值公式，
\[
\sup\{|\ell(x)|:\ x^{\mathsf T}Ax\le E\}
=
\sqrt E\,\sqrt{a^{\mathsf T}A^{-1}a}
\,=:\,R(E),
\]
从而 \(\ell(\mathcal M_A(E))\subset[-R(E),R(E)]\)。
将区间 \([-R(E),R(E)]\) 等分为 \(|\mathcal M_A(E)|-1\) 个子区间并应用鸽巢原理，存在 \(k\neq k'\) 使
\[
\abs{\ell(k)-\ell(k')}
\le
\frac{2R(E)}{|\mathcal M_A(E)|-1}.
\]
代入 \(|\mathcal M_A(E)|\) 的渐近式并整理即得
\[
\abs{\ell(k)-\ell(k')}
\le
\frac{2\sqrt{\det A}}{\operatorname{Vol}(B_r)}\sqrt{a^{\mathsf T}A^{-1}a}\,E^{\frac{1-r}{2}}
\Bigl(1+O(E^{-1/2})\Bigr),
\]
从而得到所示常数与余项阶。证毕。
\end{proof}

\begin{corollary}[单椭圆门的可辨识临界预算：由圆维控制的精度--能量反演]\label{cor:spg-single-gate-identifiability-critical-energy}
\leanverified{paper\_spg\_single\_gate\_identifiability\_critical\_energy}
在定理 \ref{thm:spg-stokes-godel-ellipsoid-min-spacing-single-gate} 的设定下，令 \(\varepsilon>0\) 为 \(\ell\) 的绝对误差容忍阈值。
若要求 \(\ell\) 在 \(\mathcal M_A(E)\) 上实现 \(\varepsilon\)-可辨识（即对任意 \(k\neq k'\) 有 \(|\ell(k)-\ell(k')|\ge \varepsilon\)），则当 \(E\to\infty\) 时必须满足
\[
E\ \le\ \left(\frac{C(a,A)}{\varepsilon}\right)^{\!\frac{2}{r-1}}\,(1+o(1)).
\]
\end{corollary}
\begin{proof}
由定理 \ref{thm:spg-stokes-godel-ellipsoid-min-spacing-single-gate}，存在 \(k\neq k'\in\mathcal M_A(E)\) 使
\(
|\ell(k)-\ell(k')|
\le
C(a,A)E^{(1-r)/2}+o(E^{(1-r)/2})
\)。
若要求所有不同对都满足 \(|\ell(k)-\ell(k')|\ge \varepsilon\)，则必有 \(C(a,A)E^{(1-r)/2}\gtrsim \varepsilon\)，等价于所示预算上界。证毕。
\end{proof}

\endinput
